\documentclass[11pt]{article}
\usepackage[margin=28mm]{geometry}
\usepackage{amsmath,amssymb,amsthm,mathtools,mathrsfs,booktabs,longtable,array}
\usepackage[colorlinks]{hyperref}
\newcommand{\statement}[3]{\subsection*{#1. #3 [#2]}}
\begin{document}
\statement{C.29}{Proved}{A full boundary test for automorphisms}
Let $S$ be a connected compact oriented surface with $m\geq1$ boundary components and negative Euler characteristic. Let $F=\widehat{\pi_1(S)}$ be its full profinite fundamental group and let $c_1,\ldots,c_m$ be its positively oriented boundary generators, based by chosen arcs. Suppose a continuous endomorphism $\varphi:F\to F$ satisfies
\[
 \varphi(c_j)=g_jc_{\sigma(j)}^{\mu_j}g_j^{-1}
 \quad(1\leq j\leq m),
 \qquad \sigma\in\mathfrak S_m,\quad
 \mu_j\in\widehat{\mathbb Z}^{\times},\quad g_j\in F.
\]
Then $\varphi$ is an automorphism. If $m\geq2$, its exponents $\mu_j$ are necessarily equal. In particular, for the full profinite pants group, a substitution
\[
 x\longmapsto x^\mu,\qquad y\longmapsto u^{-1}y^\mu u
\]
with $\mu\in\widehat{\mathbb Z}^{\times}$ is an automorphism as soon as $x^\mu u^{-1}y^\mu u$ is conjugate to $(xy)^\mu$. No separate full-generation hypothesis is needed.

\begin{proof}
Suppose the compact image $H=\varphi(F)$ is proper. Choose a proper open subgroup $U\leq F$ containing it. Writing $G=\pi_1(S)$, the subgroup $U_0=U\cap G$ has finite index $d=[F:U]>1$, and $\widehat{U_0}=U$. Let $S_U\to S$ be its connected finite covering. These are full finite-index topology comparisons, not pro-$p$ comparisons.

Replace each profinite conjugator by a discrete coset representative. Density of $G$ and openness of $U$ give $a_j\in G\cap Ug_j$. Write $a_j=v_jg_j$ with $v_j\in U$. The given image of $c_j$ is then $U$-conjugate to $a_j c_{\sigma(j)}^{\mu_j}a_j^{-1}$. As this lies in $U$ and $\mu_j$ is a full unit, closedness gives $a_j c_{\sigma(j)}a_j^{-1}\in U$. This discrete element belongs to $U_0$. The corresponding boundary component $b_j$ of $S_U$ has covering degree one over the $\sigma(j)$th base boundary, and in $H_1(U,\widehat{\mathbb Z})$ we have
\[
                       [\varphi(c_j)]=\mu_j[b_j].
\]
The selected $b_j$ are distinct because their base boundaries are distinct.

The oriented surface relation is a product of commutators times $c_1\cdots c_m=1$. Its image in the abelianization of $U$ therefore gives
\[
                         \sum_{j=1}^m\mu_j[b_j]=0.
\]
If $S_U$ has $r$ boundary components, their classes in its first integral homology have precisely one relation: their total oriented sum is zero. Equivalently, the kernel of the boundary-class map $\mathbb Z^r\to H_1(S_U,\mathbb Z)$ is generated by $(1,\ldots,1)$. This remains true over $\widehat{\mathbb Z}$: choose the usual handle and all-but-one boundary basis, so the statement follows by an explicit split free-module calculation, without an exactness assumption on an arbitrary completion. Consequently the coefficient vector having entries $\mu_j$ on the selected components and zero on every unselected component must be constant. If there were any unselected component, the constant would be zero, contradicting that every $\mu_j$ is a unit. Thus the $b_j$ are all boundary components of $S_U$.

There is therefore exactly one boundary component over each boundary of $S$, and its covering degree is one. The sum of boundary covering degrees over any one base boundary equals $d$, by the ordinary finite covering-space formula. Hence $d=1$, a contradiction. Thus $\varphi$ is surjective. A finitely generated profinite group is Hopfian, so it is an automorphism. Alternatively, preimage under a surjection permutes the finite set of open subgroups of each bounded index; every element of its kernel then lies in every open subgroup and is one.

For $m\geq2$, apply the same unique boundary relation directly in $H_1(F,\widehat{\mathbb Z})$: it gives $\mu_1=\cdots=\mu_m$. The pants specialization has $c_3=(xy)^{-1}$, and its stated third-cusp equation supplies exactly the final boundary condition. This is an independent covering proof of the pants generation step already asserted in LM4 D.35, which used the closed-subgroup rank test of LM4 D.32. It also generalizes the direct punctured-torus boundary test independently proved in C.13. No new resolution of the existence of a third-cusp parameter is claimed.
\end{proof}


\end{document}
